% GENERATED FILE. Edit canonical lessons, then run npm run generate:books.
% canonical-source-hash: fnv1a64:b374ff8cba4b025b
% canonical-lessons: MW-C174-odhno, MW-C174-odhni, MW-C174-sari, MW-C174-utaarno, MW-C174-choti

\chapter{Getting Dressed}
\label{ch:mw-getting-dressed}
\hlchaptermodality{174}

\begin{chapteropening}
\textbf{By the end of this chapter:} \emph{I can say what I put on and take off.}

Complete the last lesson of chapter 174: I can say what I put on and take off.
\end{chapteropening}

\section[\texorpdfstring{\mw{ओढणो} (oḍhṇo)}{ओढणो (oḍhṇo)}]{\mw{ओढणो} (oḍhṇo) — to cover oneself (with a cloth)}

\label{lesson:MW-C174-odhno}



\begin{quote}
Before the new one: say the Marwadi for to put out, then the Marwadi for to toast.
\end{quote}

\subsection*{You'll want to know: \mw{ओढणो}}
\textbf{\mw{ओढणो}} — \emph{oḍhṇo} — \textquotedblleft{}to cover oneself (with a cloth)\textquotedblright{}.

\begin{quote}
\mw{रात} \mw{मांय} \mw{हूं} \mw{रजाई} \mw{ओढूं} \mw{हूं।}
\end{quote}

\begin{quote}
\emph{rāt mā̃y hū̃ rajāī oḍhū̃ hū̃.}
\end{quote}

\begin{quote}
At night I pull the quilt over me.
\end{quote}

Past (I …): \mw{म्हैं} \mw{रजाई} \mw{ओढी}, \emph{mhaĩ rajāī oḍhī}.

\begin{cousinweb}
\textbf{\mw{ओढणो}} (\emph{oḍhṇo}) is inherited from a reconstructed \emph{*ōḍḍh}, \textquotedblleft{}put on, wear, cover\textquotedblright{}; CDIAL lists Old Marwari \emph{auḍhaï}, \textquotedblleft{}covers\textquotedblright{} (CDIAL 2547).
\end{cousinweb}

\begin{grammarlens}[title={What you've built}]
One more word for saying what you do.
\end{grammarlens}

\subsection*{Your turn}
\begin{itemize}
\raggedright
  \item \emph{Say it:} \emph{oḍhṇo}
  \item \emph{Say it:} \emph{oḍhṇo}, once more
  \item \emph{Say it:} \emph{rāt mā̃y hū̃ rajāī oḍhū̃ hū̃.} — At night I pull the quilt over me.
  \item \emph{Recall:} say the Marwadi for to put out, then the Marwadi for to toast, then say \emph{oḍhṇo} again
\end{itemize}

\begin{tcolorbox}[breakable,colback=teal!4,colframe=teal!35!black,title={Before you move on}]
What does \mw{ओढणो} mean? (\textquotedblleft{}To cover oneself (with a cloth)\textquotedblright{}.) Say it once more.
\end{tcolorbox}
\section[\texorpdfstring{\mw{ओढणी} (oḍhṇī)}{ओढणी (oḍhṇī)}]{\mw{ओढणी} (oḍhṇī) — a head-cloth worn by women}

\label{lesson:MW-C174-odhni}



\begin{quote}
Before the new one: say the Marwadi for to toast, then the Marwadi for to cover oneself.
\end{quote}

\subsection*{You'll want to know: \mw{ओढणी}}
\textbf{\mw{ओढणी}} — \emph{oḍhṇī} — \textquotedblleft{}a head-cloth worn by women\textquotedblright{}.

\begin{quote}
\mw{म्हारी} \mw{ओढणी} \mw{लाल} \mw{है।}
\end{quote}

\begin{quote}
\emph{mhārī oḍhṇī lāl hai.}
\end{quote}

\begin{quote}
My head-cloth is red.
\end{quote}

A feminine noun.

\begin{cousinweb}
\textbf{\mw{ओढणी}} (\emph{oḍhṇī}) is inherited from a reconstructed \emph{*ōḍḍh}, \textquotedblleft{}put on, wear, cover\textquotedblright{}, which CDIAL takes from Dravidian; it lists the Old Marwari verb \emph{auḍhaï}, \textquotedblleft{}covers\textquotedblright{}, under it (CDIAL 2547). It is the cloth you cover yourself with: \textbf{\mw{ओढणो}}.
\end{cousinweb}

\begin{grammarlens}[title={What you've built}]
One more word for saying what you do.
\end{grammarlens}

\subsection*{Your turn}
\begin{itemize}
\raggedright
  \item \emph{Say it:} \emph{oḍhṇī}
  \item \emph{Say it:} \emph{oḍhṇī}, once more
  \item \emph{Say it:} \emph{mhārī oḍhṇī lāl hai.} — My head-cloth is red.
  \item \emph{Recall:} say the Marwadi for to toast, then the Marwadi for to cover oneself, then say \emph{oḍhṇī} again
\end{itemize}

\begin{tcolorbox}[breakable,colback=teal!4,colframe=teal!35!black,title={Before you move on}]
What does \mw{ओढणी} mean? (\textquotedblleft{}A head-cloth worn by women\textquotedblright{}.) Say it once more.
\end{tcolorbox}
\section[\texorpdfstring{\mw{साड़ी} (sāṛī)}{साड़ी (sāṛī)}]{\mw{साड़ी} (sāṛī) — a sari}

\label{lesson:MW-C174-sari}



\begin{quote}
Before the new one: say the Marwadi for to cover oneself, then the Marwadi for a head-cloth.
\end{quote}

\subsection*{You'll want to know: \mw{साड़ी}}
\textbf{\mw{साड़ी}} — \emph{sāṛī} — \textquotedblleft{}a sari\textquotedblright{}.

\begin{quote}
\mw{मां} \mw{नुवी} \mw{साड़ी} \mw{पैरै} \mw{है।}
\end{quote}

\begin{quote}
\emph{mā̃ nuvī sāṛī pairai hai.}
\end{quote}

\begin{quote}
Mother wears a new sari.
\end{quote}

A feminine noun.

\begin{cousinweb}
\textbf{\mw{साड़ी}} (\emph{sāṛī}) is inherited from Sanskrit \emph{śāṭī}, \textquotedblleft{}a strip of cloth\textquotedblright{}; CDIAL lists Old Marwari \emph{sāḍī}, \textquotedblleft{}a sari\textquotedblright{}, feminine (CDIAL 12381).
\end{cousinweb}

\begin{grammarlens}[title={What you've built}]
One more word for saying what you do.
\end{grammarlens}

\subsection*{Your turn}
\begin{itemize}
\raggedright
  \item \emph{Say it:} \emph{sāṛī}
  \item \emph{Say it:} \emph{sāṛī}, once more
  \item \emph{Say it:} \emph{mā̃ nuvī sāṛī pairai hai.} — Mother wears a new sari.
  \item \emph{Recall:} say the Marwadi for to cover oneself, then the Marwadi for a head-cloth, then say \emph{sāṛī} again
\end{itemize}

\begin{tcolorbox}[breakable,colback=teal!4,colframe=teal!35!black,title={Before you move on}]
What does \mw{साड़ी} mean? (\textquotedblleft{}A sari\textquotedblright{}.) Say it once more.
\end{tcolorbox}
\section[\texorpdfstring{\mw{उतारणो} (utārṇo)}{उतारणो (utārṇo)}]{\mw{उतारणो} (utārṇo) — to take off (clothes)}

\label{lesson:MW-C174-utaarno}



\begin{quote}
Before the new one: say the Marwadi for a head-cloth, then the Marwadi for a sari.
\end{quote}

\subsection*{You'll want to know: \mw{उतारणो}}
\textbf{\mw{उतारणो}} — \emph{utārṇo} — \textquotedblleft{}to take off (clothes)\textquotedblright{}.

\begin{quote}
\mw{हूं} \mw{घरै} \mw{जूती} \mw{उतारूं} \mw{हूं।}
\end{quote}

\begin{quote}
\emph{hū̃ gharai jūtī utārū̃ hū̃.}
\end{quote}

\begin{quote}
At home I take off my shoes.
\end{quote}

Past (I …): \mw{म्हैं} \mw{जूती} \mw{उतारी}, \emph{mhaĩ jūtī utārī}.

\begin{cousinweb}
\textbf{\mw{उतारणो}} (\emph{utārṇo}) is inherited from Sanskrit \emph{uttārayati}, the causative of \emph{uttarati}, \textquotedblleft{}comes down\textquotedblright{}; CDIAL lists Old Marwari \emph{utaraï}, \textquotedblleft{}descends\textquotedblright{} (CDIAL 1770). Its partner \textbf{\mw{उतरणो}}, to get down, is already yours.
\end{cousinweb}

\begin{grammarlens}[title={What you've built}]
One more word for saying what you do.
\end{grammarlens}

\subsection*{Your turn}
\begin{itemize}
\raggedright
  \item \emph{Say it:} \emph{utārṇo}
  \item \emph{Say it:} \emph{utārṇo}, once more
  \item \emph{Say it:} \emph{hū̃ gharai jūtī utārū̃ hū̃.} — At home I take off my shoes.
  \item \emph{Recall:} say the Marwadi for a head-cloth, then the Marwadi for a sari, then say \emph{utārṇo} again
\end{itemize}

\begin{tcolorbox}[breakable,colback=teal!4,colframe=teal!35!black,title={Before you move on}]
What does \mw{उतारणो} mean? (\textquotedblleft{}To take off (clothes)\textquotedblright{}.) Say it once more.
\end{tcolorbox}
\section[\texorpdfstring{\mw{चोटी} (choṭī)}{चोटी (choṭī)}]{\mw{चोटी} (choṭī) — a plait of hair}

\label{lesson:MW-C174-choti}



\begin{quote}
Before the new one: say the Marwadi for a sari, then the Marwadi for to take off.
\end{quote}

\subsection*{You'll want to know: \mw{चोटी}}
\textbf{\mw{चोटी}} — \emph{choṭī} — \textquotedblleft{}a plait of hair\textquotedblright{}.

\begin{quote}
\mw{म्हारी} \mw{चोटी} \mw{लांबी} \mw{है।}
\end{quote}

\begin{quote}
\emph{mhārī choṭī lā̃bī hai.}
\end{quote}

\begin{quote}
My plait is long.
\end{quote}

A feminine noun.

\begin{cousinweb}
\textbf{\mw{चोटी}} (\emph{choṭī}) is inherited from Sanskrit \emph{cūḍa}, \textquotedblleft{}a topknot\textquotedblright{}; CDIAL lists Gujarati \emph{coṭī}, \textquotedblleft{}a tuft of hair\textquotedblright{} (CDIAL 4883). CDIAL gives no Marwari form in this entry.
\end{cousinweb}

\begin{grammarlens}[title={What you've built}]
One more word for saying what you do.
\end{grammarlens}

\subsection*{Your turn}
\begin{itemize}
\raggedright
  \item \emph{Say it:} \emph{choṭī}
  \item \emph{Say it:} \emph{choṭī}, once more
  \item \emph{Say it:} \emph{mhārī choṭī lā̃bī hai.} — My plait is long.
  \item \emph{Recall:} say the Marwadi for a sari, then the Marwadi for to take off, then say \emph{choṭī} again
\end{itemize}

\begin{tcolorbox}[breakable,colback=teal!4,colframe=teal!35!black,title={Before you move on}]
What does \mw{चोटी} mean? (\textquotedblleft{}A plait of hair\textquotedblright{}.) Say it once more.
\end{tcolorbox}
